Bei der snoopmedia GmbH handelt es sich um eine inhabergeführte Digitalagentur mit Sitz in der Otto-Hahn-Straße in Grafschaft bei Bonn. Seit ihrer Gründung realisiert das Unternehmen 
technologisch anspruchsvolle Softwarelösungen und hat sich als Spezialist für die Digitalisierung komplexer Geschäftsprozesse etabliert. Das Leistungsspektrum ist breit gefächert und 
umfasst neben der klassischen Web- und App-Entwicklung auch die Betreuung umfangreicher E-Commerce-Systeme sowie die Durchführung innovativer Forschungsprojekte in den Bereichen 
Künstliche Intelligenz (KI) und Internet of Things (IoT). Besonders zeichnet sich snoopmedia durch die Beteiligung an öffentlich geförderten Projekten wie „VERO“, „CitiDigiSpace“ 
oder aktuell „GameUp“ aus, bei denen zukunftsweisende digitale Ökosysteme in Zusammenarbeit mit akademischen Partnern entwickelt werden.

Die technologische Basis der Agentur ist konsequent auf Open-Source-Technologien ausgerichtet. Hierbei spielen Frameworks wie Symfony (PHP) sowie moderne JavaScript-Bibliotheken 
und MySQL-Datenbanken eine zentrale Rolle. Intern arbeitet das Team nach agilen Prinzipien, was eine enge Verzahnung zwischen Konzeption und technischer Umsetzung ermöglicht. Für 
mein Praxisprojekt bietet diese Struktur den Vorteil, dass ich in ein erfahrenes Team von Softwareentwicklern integriert bin, die sowohl kundenorientierte Lösungen als auch interne 
technologische Innovationen vorantreiben. Die snoopmedia GmbH erfüllt somit nicht nur die formalen Anforderungen an eine fachkundige Betreuung durch erfahrene Ansprechpartner, sondern 
bietet durch ihre Spezialisierung auf individuelle Softwarelösungen genau das Umfeld für die von mir angestrebte vertiefte Programmiertätigkeit.